\section{2024-10-23}

pascalsches dreieck

\startframedtipp

\definemathmatrix[pmatrix]        % defining matrix with parentheses
	[matrix:parentheses]
	[simplecommand=pmatrix]

\startformula
\startpmatrix
	   \NC 5 \NR
	   \NC 2 \NR
\stoppmatrix = 10
\stopformula

Fünf über drei

Heißt drei von 5 Treffer

Binominalkoefficient

\stopframedtipp


\startbuffer[mp:fig:unterricht:binominal:pascal]
\startluacode
  local function fact (n)
    if n <= 0 then
      return 1
    else
      return n * fact(n-1)
    end
  end

  local function ncr(n,r)
    return fact(n)/(fact(r)*fact(n-r))
  end
  userdata.P = {
       fact = fact,
       ncr  = ncr,
  }
  function MP.p_ncr(n, r)
    mp.print(ncr(n,r))
  end
  function MP.p_ncr_x()
    mp.print(ncr(mp.scan.pair()))
  end
\stopluacode

\startMPpage

numeric n,r,u,dx,dy,tt,final; u := 0.8cm;

path p, q;
pair A,now;

A := dir(225)*u;
dy := ypart A;
dx := -2*(xpart A);

for n=0 upto 8:
  for r=0 upto n:
%   tt := lua("mp.print(P.ncr(" & decimal n & "," &  decimal r & " ))");
  % tt := lua.MP.p_ncr(n,r);
    tt := runscript("MP.p_ncr_x()") (n,r)  ;
    now := n*A + (r*dx,0);
    label(textext("$" & decimal tt & "$"),now);
   endfor;
endfor;

\stopMPpage
\stopbuffer

\placefigure[force][mp:fig:unterricht:binominal:pascal]{Pascal's dreieck}{
\typesetbuffer[mp:fig:unterricht:binominal:pascal]
}
